
\documentclass[12pt,a4paper]{article}

% ============================================================
% PACOTES
% ============================================================

\usepackage[utf8]{inputenc}
\usepackage[T1]{fontenc}
\usepackage[brazil]{babel}
\usepackage{lmodern}
\usepackage{geometry}
\usepackage{setspace}
\usepackage{booktabs}
\usepackage{array}
\usepackage{amsmath}
\usepackage{amssymb}
\usepackage{graphicx}
\usepackage{float}
\usepackage{hyperref}
\usepackage{xcolor}
\usepackage{longtable}

\geometry{
    left=2.5cm,
    right=2.5cm,
    top=2.5cm,
    bottom=2.5cm
}

\onehalfspacing

\hypersetup{
    colorlinks=true,
    linkcolor=blue,
    urlcolor=blue,
    citecolor=blue
}

% ============================================================
% DOCUMENTO
% ============================================================

\begin{document}

\begin{titlepage}
    \centering

    \vspace*{2cm}

    {\Huge\bfseries IONA\par}

    \vspace{0.5cm}

    {\Large Earth Intelligence \& Pattern Discovery\par}

    \vspace{1.5cm}

    {\LARGE\bfseries Revealing What the Averages Hide\par}

    \vspace{2cm}

    {\Large Relatório Técnico-Científico\par}

    \vspace{1cm}

    {\large Análise espaço-temporal da precipitação na
    região-piloto do Cerrado brasileiro\par}

    \vfill

    {\large NASA Space Apps Challenge 2026\par}

    \vspace{0.5cm}

    {\large 14 de novembro de 2026\par}

\end{titlepage}

\tableofcontents
\newpage

% ============================================================
\section{Resumo}
% ============================================================

O IONA (Earth Intelligence \& Pattern Discovery) é uma
plataforma de análise e descoberta de padrões ambientais
que investiga estruturas espaço-temporais que podem não
ser evidentes em estatísticas regionais agregadas.

Este estudo investiga a seguinte questão:

\begin{quote}
\textbf{A precipitação média do Cerrado brasileiro oculta
tendências regionais divergentes?}
\end{quote}

Como estudo-piloto, foi selecionada uma área retangular
nas proximidades do Distrito Federal e regiões adjacentes.
Foram utilizados dados mensais do produto NASA GPM IMERG
Monthly Final Run, versão 07, abrangendo janeiro de 2001
a dezembro de 2024.

A série foi convertida de taxa média de precipitação para
acumulação mensal, permitindo a construção de um conjunto
espaço-temporal com 288 meses e 625 células espaciais.

A análise incluiu estimativas de tendência de Sen,
teste de Mann--Kendall, correção de múltiplas comparações
por Benjamini--Hochberg, diagnóstico de autocorrelação,
Moran Global e Moran Local (LISA).

Os resultados indicaram forte autocorrelação espacial
nas estimativas de tendência. Entretanto, não foram
identificadas tendências temporais estatisticamente
significativas após a correção FDR nas células analisadas.
A análise LISA revelou agrupamentos espaciais de
magnitudes relativamente altas e baixas, mas sua
delimitação apresentou sensibilidade à matriz de
vizinhança utilizada.

O estudo demonstra a importância de investigar a
distribuição espacial dos fenômenos ambientais, sem
confundir organização espacial das estimativas com
significância estatística de tendências temporais.

% ============================================================
\section{Pergunta de pesquisa e objetivos}
% ============================================================

\subsection{Pergunta de pesquisa}

A precipitação média regional oculta diferenças
espaciais nas magnitudes e direções das tendências
temporais de precipitação?

\subsection{Objetivo geral}

Investigar padrões espaço-temporais de precipitação
em uma região-piloto do Cerrado brasileiro utilizando
dados de sensoriamento remoto e métodos estatísticos
espaciais.

\subsection{Objetivos específicos}

\begin{itemize}
    \item Construir uma série mensal de precipitação
    a partir dos dados NASA GPM IMERG.
    \item Caracterizar a sazonalidade e a variabilidade
    interanual da precipitação.
    \item Estimar tendências temporais regionais e locais.
    \item Investigar autocorrelação espacial nas
    estimativas de tendência.
    \item Identificar agrupamentos espaciais de
    magnitudes semelhantes.
    \item Avaliar a sensibilidade dos resultados à
    definição de vizinhança espacial.
\end{itemize}

% ============================================================
\section{Área de estudo}
% ============================================================

A área de estudo corresponde a um retângulo geográfico
utilizado como região-piloto, delimitado pelas seguintes
coordenadas:

\begin{itemize}
    \item Longitude oeste: $-48.5^\circ$
    \item Longitude leste: $-46.0^\circ$
    \item Latitude sul: $-16.5^\circ$
    \item Latitude norte: $-14.0^\circ$
\end{itemize}

A grade espacial possui resolução de aproximadamente
$0.1^\circ$, resultando em 25 posições de longitude e
25 posições de latitude, totalizando 625 células.

\textbf{Delimitação importante:} esse retângulo é uma
área-piloto de análise e não representa o polígono
oficial completo do bioma Cerrado. Portanto, os
resultados não devem ser generalizados automaticamente
para todo o bioma.

% ============================================================
\section{Fonte e características dos dados}
% ============================================================

\subsection{NASA GPM IMERG}

Foi utilizado o produto:

\begin{itemize}
    \item Missão: Global Precipitation Measurement (GPM).
    \item Produto: Integrated Multi-satellitE Retrievals
    for GPM (IMERG).
    \item Produto mensal: 3IMERGM.
    \item Versão: V07 / V07B.
    \item Identificador: \texttt{GPM\_3IMERGM}.
    \item Período: janeiro de 2001 a dezembro de 2024.
    \item Frequência: mensal.
    \item Resolução espacial: aproximadamente $0.1^\circ$.
\end{itemize}

Identificador DOI:

\begin{center}
\url{https://doi.org/10.5067/GPM/IMERG/3B-MONTH/07}
\end{center}

Página oficial do produto:

\begin{center}
\url{https://disc.gsfc.nasa.gov/datasets/GPM_3IMERGM_07/summary}
\end{center}

\subsection{Conversão da variável}

O produto mensal fornece a taxa média de precipitação
ao longo do mês, em milímetros por hora (mm/hr).

Para obter a acumulação mensal em milímetros, foi
utilizada a conversão:

\begin{equation}
P_m = R_m \times 24 \times D_m
\end{equation}

em que:

\begin{itemize}
    \item $P_m$ é a precipitação acumulada no mês,
    em mm/mês;
    \item $R_m$ é a taxa média de precipitação,
    em mm/hr;
    \item $D_m$ é o número de dias do mês.
\end{itemize}

A duração real de cada mês foi considerada,
incluindo anos bissextos.

Os valores de preenchimento (\textit{fill values})
foram convertidos para valores ausentes (NaN).

\subsection{Validação}

Foram encontrados 288 arquivos mensais, correspondentes
aos 288 meses esperados no período estudado.

\begin{table}[H]
\centering
\caption{Validação dos dados de entrada}
\begin{tabular}{lr}
\toprule
Indicador & Resultado \\
\midrule
Arquivos esperados & 288 \\
Arquivos encontrados & 288 \\
Arquivos HDF5 válidos & 288 \\
Meses ausentes & 0 \\
Arquivos duplicados & 0 \\
Arquivos inválidos & 0 \\
Cobertura espacial média & 100\% \\
Dados ausentes no cubo final & 0\% \\
\bottomrule
\end{tabular}
\end{table}

O conjunto final possui dimensão:

\begin{equation}
288 \times 25 \times 25
\end{equation}

correspondente a 288 meses e 625 células espaciais.

% ============================================================
\section{Metodologia}
% ============================================================

\subsection{Caracterização exploratória}

Foi calculada a climatologia mensal considerando os
24 anos disponíveis.

Também foi construída uma série anual regional,
obtida pela soma das precipitações mensais após a
agregação espacial ponderada.

A ponderação espacial utiliza:

\begin{equation}
w_i = \cos(\phi_i)
\end{equation}

em que $\phi_i$ representa a latitude da célula.
Essa aproximação considera a variação da área
representada pelas células em diferentes latitudes.

A média regional ponderada é:

\begin{equation}
\bar{P}_t =
\frac{\sum_i w_i P_{i,t}}
{\sum_i w_i}
\end{equation}

\subsection{Tendência temporal}

A magnitude das tendências foi estimada pelo
estimador de Sen:

\begin{equation}
\beta =
\operatorname{mediana}
\left(
\frac{x_j-x_i}{j-i}
\right),
\qquad j>i
\end{equation}

O teste de Mann--Kendall foi utilizado para avaliar
a presença de tendência monotônica.

Foram calculados intervalos de confiança de 95\%
para as estimativas de Sen.

As análises locais foram realizadas individualmente
para as 625 células espaciais.

\subsection{Correção de múltiplas comparações}

Como foram realizados testes simultâneos em diversas
células, foi aplicada a correção de Benjamini--Hochberg
(BH), controlando a taxa de falsas descobertas (FDR).

A significância estatística local foi avaliada
considerando:

\begin{equation}
q_i < 0.05
\end{equation}

A correção foi aplicada separadamente aos resultados
convencionais de Mann--Kendall e à análise de
sensibilidade com TFPW.

\subsection{Autocorrelação temporal}

A autocorrelação de defasagem 1 foi examinada
tanto na série original quanto nos resíduos após
a remoção da tendência estimada.

Foi realizada uma análise de sensibilidade utilizando
Trend-Free Pre-Whitening (TFPW).

Essa análise foi tratada como diagnóstico complementar,
não como correção definitiva da inferência estatística.

\subsection{Autocorrelação espacial}

A dependência espacial das estimativas de Sen foi
investigada utilizando o Moran Global:

\begin{equation}
I =
\frac{n}{S_0}
\frac{
\sum_i \sum_j w_{ij} z_i z_j
}{
\sum_i z_i^2
}
\end{equation}

em que:

\begin{itemize}
    \item $n$ é o número de células;
    \item $w_{ij}$ representa a relação de vizinhança;
    \item $z_i$ representa o desvio da estimativa em
    relação à média.
\end{itemize}

Foram utilizadas 9.999 permutações.

\subsection{Moran Local (LISA)}

O Moran Local foi utilizado para identificar
agrupamentos espaciais:

\begin{itemize}
    \item HH: valores altos próximos de valores altos;
    \item LL: valores baixos próximos de valores baixos;
    \item HL: valores altos próximos de valores baixos;
    \item LH: valores baixos próximos de valores altos.
\end{itemize}

Os testes locais também foram submetidos à correção
FDR por Benjamini--Hochberg.

Os quadrantes representam relações espaciais
relativas à distribuição das estimativas de Sen.
Não representam, por si só, tendências temporais
positivas ou negativas estatisticamente significativas.

\subsection{Índice de cancelamento espacial do IONA}

Foi definido um indicador descritivo próprio,
denominado Spatial Trend Cancellation Index (STCI):

\begin{equation}
STCI =
1 -
\frac{
\left|\sum_i w_i \beta_i\right|
}{
\sum_i w_i |\beta_i|
}
\end{equation}

em que $\beta_i$ representa a estimativa de Sen
para a célula $i$ e $w_i$ representa seu peso espacial.

O índice varia entre zero e um:

\begin{itemize}
    \item Valores próximos de zero indicam pouca
    compensação entre magnitudes positivas e negativas.
    \item Valores próximos de um indicam maior
    compensação entre contribuições positivas e negativas.
\end{itemize}

O STCI é uma métrica descritiva definida pelo projeto
IONA, não um índice climático padronizado.

O indicador utiliza estimativas pontuais e não
incorpora a incerteza estatística das inclinações.

\subsection{Robustez da vizinhança espacial}

A estabilidade dos resultados espaciais foi avaliada
comparando duas matrizes de contiguidade:

\begin{itemize}
    \item Queen: considera vizinhos laterais e diagonais.
    \item Rook: considera somente vizinhos que compartilham
    uma aresta.
\end{itemize}

A concordância das classificações foi calculada
diretamente. A sobreposição dos agrupamentos foi
avaliada pelo índice de Jaccard:

\begin{equation}
J(A,B) =
\frac{|A \cap B|}{|A \cup B|}
\end{equation}

% ============================================================
\section{Resultados}
% ============================================================

\subsection{Sazonalidade}

A climatologia mensal apresentou forte sazonalidade,
com maiores acumulados entre outubro e abril e
menores valores entre maio e setembro.

\begin{table}[H]
\centering
\caption{Climatologia mensal regional}
\begin{tabular}{lc}
\toprule
Mês & Precipitação média (mm) \\
\midrule
Janeiro & 221.45 \\
Fevereiro & 182.95 \\
Março & 200.35 \\
Abril & 97.18 \\
Maio & 18.60 \\
Junho & 2.49 \\
Julho & 0.56 \\
Agosto & 6.76 \\
Setembro & 24.33 \\
Outubro & 109.18 \\
Novembro & 230.44 \\
Dezembro & 243.03 \\
\bottomrule
\end{tabular}
\end{table}

\subsection{Tendência regional}

A estimativa regional de Sen foi:

\begin{equation}
\beta = -11.048
\end{equation}

em unidades de mm/ano de precipitação anual
por ano de calendário.

O intervalo de confiança de 95\% foi:

\begin{equation}
IC_{95\%} = [-27.679,\ 5.828]
\end{equation}

O teste de Mann--Kendall apresentou:

\begin{itemize}
    \item Kendall's tau: $-0.2174$;
    \item valor-p: $0.144623$.
\end{itemize}

O intervalo de confiança inclui zero e o valor-p
é superior a 0.05. Portanto, não há evidência
estatística suficiente, por esses testes, para
afirmar uma tendência monotônica regional no
período analisado.

\subsection{Tendências locais}

Foram avaliadas 625 células.

\begin{table}[H]
\centering
\caption{Resultados dos testes temporais locais}
\begin{tabular}{lr}
\toprule
Indicador & Resultado \\
\midrule
Células analisadas & 625 \\
Significativas após FDR (MK) & 0 \\
Significativas após FDR (TFPW) & 0 \\
Significativas em ambos & 0 \\
Células classificadas como aumento robusto & 0 \\
Células classificadas como redução robusta & 0 \\
\bottomrule
\end{tabular}
\end{table}

A ausência de células significativas após FDR não
demonstra que todas as inclinações sejam nulas.
Indica que os testes realizados não forneceram
evidência suficiente para declarar tendências
temporais locais significativas após a correção
de múltiplas comparações.

\subsection{Distribuição espacial das inclinações}

A distribuição das estimativas de Sen apresentou:

\begin{itemize}
    \item 618 células com inclinação negativa;
    \item 7 células com inclinação positiva;
    \item média das inclinações: $-10.4161$;
    \item mediana: $-10.6109$;
    \item mínimo: $-16.7726$;
    \item máximo: $+6.3680$.
\end{itemize}

As sete estimativas positivas se concentraram no
setor noroeste da região-piloto.

Entretanto, nenhuma dessas células apresentou
significância temporal após FDR.

\subsection{Spatial Trend Cancellation Index}

O STCI calculado foi:

\begin{equation}
STCI = 0.0056
\end{equation}

Equivalente a aproximadamente 0.56\%.

As contribuições ponderadas foram:

\begin{itemize}
    \item contribuição positiva: aproximadamente 0.28\%;
    \item contribuição negativa: aproximadamente 99.72\%.
\end{itemize}

O resultado indica compensação muito pequena entre
as inclinações positivas e negativas estimadas.

Portanto, neste recorte e período, o índice não
aponta cancelamento expressivo entre tendências
de direções opostas.

Esse resultado é descritivo e não deve ser
interpretado como confirmação de uma tendência
climática regional.

\subsection{Moran Global}

A análise de autocorrelação espacial utilizando
contiguidade Queen apresentou:

\begin{equation}
I = 0.852985
\end{equation}

com valor-p por permutação igual a 0.0001.

O resultado indica forte autocorrelação espacial
positiva nas estimativas de Sen: células próximas
tendem a apresentar magnitudes semelhantes.

Esse teste avalia a organização espacial das
estimativas, não a significância temporal das
tendências.

\subsection{Moran Local (LISA)}

A análise Queen identificou 322 células significativas
após FDR:

\begin{table}[H]
\centering
\caption{Agrupamentos LISA com contiguidade Queen}
\begin{tabular}{lr}
\toprule
Classificação & Células \\
\midrule
High-High & 136 \\
Low-Low & 186 \\
High-Low & 0 \\
Low-High & 0 \\
Não significativas & 303 \\
\bottomrule
\end{tabular}
\end{table}

Os agrupamentos HH e LL representam regiões com
magnitudes relativamente altas e baixas das
estimativas de Sen, respectivamente.

Importante: HH não significa necessariamente
tendência temporal positiva, assim como LL não
significa necessariamente tendência temporal negativa.

A maioria das células HH também apresenta
inclinações negativas. Portanto, os agrupamentos
separam principalmente magnitudes relativas,
e não direções temporais opostas.

\subsection{Contraste HH versus LL}

A comparação exploratória apresentou:

\begin{table}[H]
\centering
\caption{Contraste das estimativas de Sen}
\begin{tabular}{lrr}
\toprule
Estatística & HH & LL \\
\midrule
Número de células & 136 & 186 \\
Média & -5.9633 & -13.8966 \\
Mediana & -6.4655 & -13.9218 \\
\bottomrule
\end{tabular}
\end{table}

A diferença entre as médias foi:

\begin{equation}
\Delta_{\mathrm{HH-LL}} = 7.9332
\end{equation}

O teste exploratório de permutação apresentou
valor-p de 0.0001.

Esse teste não considera a dependência espacial
nem o fato de os grupos terem sido definidos
a partir das próprias estimativas analisadas.
Consequentemente, seu valor-p não deve ser
interpretado como evidência inferencial independente.

O contraste descreve uma diferença de magnitude
entre os agrupamentos identificados.

\subsection{Robustez Queen versus Rook}

A comparação das matrizes de vizinhança produziu:

\begin{table}[H]
\centering
\caption{Comparação das matrizes de contiguidade}
\begin{tabular}{lrr}
\toprule
Indicador & Queen & Rook \\
\midrule
Moran Global & 0.852985 & 0.883403 \\
Valor-p global & 0.0001 & 0.0001 \\
Células LISA significativas & 322 & 139 \\
HH & 136 & 43 \\
LL & 186 & 96 \\
\bottomrule
\end{tabular}
\end{table}

Indicadores de concordância:

\begin{itemize}
    \item Concordância das classificações: 70.72\%;
    \item Jaccard dos agrupamentos significativos: 0.432;
    \item Jaccard HH: 0.316;
    \item Jaccard LL: 0.516.
\end{itemize}

O Moran Global permaneceu elevado nas duas
configurações, indicando que a autocorrelação
espacial geral é consistente.

Entretanto, a delimitação dos agrupamentos locais
apresentou sensibilidade à definição de vizinhança.

Assim, os agrupamentos LISA devem ser apresentados
como resultados exploratórios, e não como
fronteiras espaciais definitivas.

% ============================================================
\section{Discussão}
% ============================================================

Os resultados demonstram que a análise espacial
pode revelar estruturas que não são descritas
adequadamente por uma única estatística regional.

A forte autocorrelação espacial das inclinações
indica que as estimativas apresentam organização
geográfica: células próximas tendem a possuir
magnitudes semelhantes.

Entretanto, a ausência de significância temporal
após FDR impede afirmar que essas inclinações
representem tendências temporais estatisticamente
confirmadas.

A análise LISA identificou agrupamentos de
magnitudes relativamente altas e baixas. A
comparação entre HH e LL indicou uma diferença
descritiva relevante entre os valores estimados.

Contudo, a sensibilidade à matriz de vizinhança
mostra que a delimitação desses agrupamentos
não é inteiramente estável.

O STCI também apresentou baixo cancelamento
entre contribuições positivas e negativas.
Portanto, os resultados deste recorte não
sustentam a hipótese de que uma média regional
esteja escondendo uma forte compensação entre
tendências temporais de sinais opostos.

Essas conclusões são específicas da região-piloto,
do produto utilizado, da resolução espacial e
do período estudado.

% ============================================================
\section{Limitações}
% ============================================================

As principais limitações são:

\begin{enumerate}
    \item \textbf{Recorte geográfico:} a área utilizada
    é um retângulo-piloto e não o polígono oficial
    completo do Cerrado.

    \item \textbf{Período:} 24 anos representam uma
    série relativamente curta para caracterização
    de tendências climáticas de longo prazo.

    \item \textbf{Autocorrelação temporal:} os testes
    convencionais podem ser afetados pela dependência
    serial. A análise TFPW utilizada é uma análise
    de sensibilidade e não resolve integralmente
    a incerteza associada à estimação da tendência
    e da autocorrelação.

    \item \textbf{Dependência espacial:} a correção
    FDR convencional não elimina a dependência
    entre células vizinhas.

    \item \textbf{Sensibilidade LISA:} os agrupamentos
    locais variam conforme a matriz de vizinhança.

    \item \textbf{Incerteza do STCI:} o índice utiliza
    estimativas pontuais e não incorpora intervalos
    de confiança ou propagação de incerteza.

    \item \textbf{Dados de precipitação:} produtos
    de sensoriamento remoto possuem incertezas
    próprias de estimativa e recuperação.
\end{enumerate}

% ============================================================
\section{Conclusões}
% ============================================================

A análise realizada pelo IONA identificou forte
organização espacial nas estimativas de tendência
da precipitação na região-piloto.

O Moran Global apresentou valores elevados tanto
para a matriz Queen quanto para a matriz Rook,
indicando que a autocorrelação espacial geral
não depende exclusivamente da inclusão de
vizinhos diagonais.

Por outro lado, a classificação local LISA
apresentou sensibilidade à matriz de vizinhança,
exigindo cautela na interpretação dos agrupamentos.

Os testes temporais não identificaram tendências
regionais ou locais estatisticamente significativas
após os procedimentos de avaliação e correção
utilizados.

O STCI calculado também indicou baixo cancelamento
entre as estimativas positivas e negativas.

Portanto, os resultados deste estudo-piloto não
confirmam a hipótese de forte cancelamento de
tendências temporais opostas na região analisada.

O principal resultado metodológico do IONA é
demonstrar a importância de separar três dimensões:

\begin{enumerate}
    \item magnitude e direção das estimativas
    temporais;
    \item significância estatística das tendências;
    \item organização espacial das estimativas.
\end{enumerate}

Essa separação permite produzir interpretações
mais rigorosas e evita que agrupamentos espaciais
sejam confundidos com evidências de mudança
climática temporal.

% ============================================================
\section{Trabalhos futuros}
% ============================================================

As próximas etapas de desenvolvimento do IONA incluem:

\begin{itemize}
    \item Aplicar uma máscara geográfica oficial
    do bioma Cerrado.
    \item Investigar a estabilidade das estimativas
    com métodos adicionais para autocorrelação temporal.
    \item Avaliar incertezas associadas às estimativas
    de tendência e ao STCI.
    \item Integrar outras variáveis ambientais
    disponibilizadas por fontes NASA.
    \item Desenvolver detecção de anomalias e
    identificação de padrões com aprendizado de máquina.
    \item Construir uma interface interativa para
    exploração dos resultados.
    \item Implementar geração de relatórios com
    rastreabilidade dos dados e métodos.
\end{itemize}

% ============================================================
\section*{Referências}
% ============================================================

\begin{enumerate}

    \item NASA. \textit{Global Precipitation Measurement
    (GPM) -- IMERG}. Disponível em:
    \url{https://gpm.nasa.gov/data/imerg}.

    \item NASA GES DISC. \textit{GPM IMERG Monthly Final
    Run V07}. DOI:
    \url{https://doi.org/10.5067/GPM/IMERG/3B-MONTH/07}.

    \item MANN, H. B. Nonparametric Tests Against Trend.
    \textit{Econometrica}, v. 13, n. 3, p. 245--259, 1945.

    \item KENDALL, M. G. \textit{Rank Correlation Methods}.
    London: Griffin, 1975.

    \item SEN, P. K. Estimates of the Regression
    Coefficient Based on Kendall's Tau.
    \textit{Journal of the American Statistical
    Association}, v. 63, n. 324, p. 1379--1389, 1968.

    \item BENJAMINI, Y.; HOCHBERG, Y. Controlling the
    False Discovery Rate: A Practical and Powerful
    Approach to Multiple Testing.
    \textit{Journal of the Royal Statistical Society:
    Series B}, v. 57, n. 1, p. 289--300, 1995.

    \item MORAN, P. A. P. Notes on Continuous Stochastic
    Phenomena. \textit{Biometrika}, v. 37, n. 1/2,
    p. 17--23, 1950.

    \item ANSELIN, L. Local Indicators of Spatial
    Association---LISA.
    \textit{Geographical Analysis}, v. 27, n. 2,
    p. 93--115, 1995.

\end{enumerate}

\end{document}
